\section{Introducción}

En un mercado retail y comercial altamente competitivo como el de Arequipa, la capacidad de respuesta rápida y la gestión eficiente de la información logística son pilares fundamentales para la sostenibilidad y el crecimiento de las pequeñas y medianas empresas. En este contexto, Industria Textil Atlas E.I.R.L., liderada por su gerente general Víctor Andre Huayllaro Guevara, ha consolidado su presencia en el sector deportivo a través de la venta y distribución de artículos de protección y accesorios deportivos (rodilleras, musleras, tobilleras, coderas, pelotas, bandas de resistencia, etc.) comercializados desde su local físico en San Camilo y de forma digital.

No obstante, a pesar de su sólida trayectoria de más de 20 años, la empresa operaba bajo un modelo de madurez digital inicial (estimado en 35.34\%), caracterizado por una marcada dependencia de procesos manuales. Al no contar con producción propia y tercerizar por completo la confección de sus artículos textiles a talleres externos, el flujo de trabajo diario de stock, compras y ventas en mostrador se encontraba fragmentado. La gestión dependía de cuadernos físicos, anotaciones en hojas sueltas e inspecciones visuales en el almacén, lo que dilataba los tiempos de despacho a 15--20 minutos por cliente en mostrador y generaba cuellos de botella para cotizar pedidos mayoristas.

Ante este escenario, y bajo el marco del contrato N° 1036-PROINNOVATE-IMTEMD-2025, se planteó la implementación de un sistema ERP simplificado basado en Odoo Online (SaaS) como solución integral. Bajo un enfoque puramente Out-of-the-box (nativo), el proyecto se centró en la parametrización de los módulos de Inventario, Punto de Venta (POS), Ventas y Compras. Esta transformación tecnológica buscó erradicar la fragmentación de la información de almacén con códigos SKU, automatizar las sugerencias de compra a confeccionistas externos mediante puntos de reorden, y agilizar las transacciones en tienda con la terminal POS.

El presente informe tiene como finalidad evaluar los resultados obtenidos tras la puesta en marcha de Odoo, analizando el impacto generado en la productividad operativa, la optimización de los tiempos de atención, la reducción de errores en existencias y el nivel de adopción del sistema por parte del personal, evidenciando los beneficios tangibles alcanzados para la escalabilidad del negocio.
